\chapter{What remains open}\label{ch:open}

A reader who has followed the book this far may hold a belief that the earlier
chapters have encouraged. The method has propositions with full proofs, a pipeline
that regenerates every number from one configuration, falsification tests that can
refuse a decision, a protocol with roles and a filled record. On this view the
scientific question is settled, and what remains is adoption: an institution has to
install the software, train its staff and fill in the worksheets. The view is
understandable, because the parts of the method fit together and each part has been
checked.

The view fails on one table. Every floor this book reports for the credit register
depends on how precisely the adversary can read a loan's exposure, and no institution
has measured that precision. When the standard deviation of the reading on the log
scale is 0.6, the programme floor against the full outside adversary has a median of
150 over twenty draws of the synthetic register. When it is 0.3, the median is 300,
and in three of the twenty draws no floor up to 300 passes at all. When it is 1, the
median is 75. % source: paper/generated/table_replication.tex
The decision rule is settled. The number it produces for a real register is not,
because the inputs that decide it are declarations about the world outside the
institution, and the evidence for them has not yet been gathered on any real
register.

This chapter asks which of the book's claims would change if the method were run on
a real register by a real institution, and which claims the method cannot make at
all. It closes by answering the question that Chapter~\ref{ch:decision} put at the
start of the book.

\section{Real institutional registers}

The credit register of Chapter~\ref{ch:credit} is synthetic. It shares the structure
of a supervisory collection such as AnaCredit \citep{israel2017,ecb2016}: firms as
borrowers, lenders with unequal shares, sector, region and size class as the key,
monthly vintages. Its parameters were chosen, not estimated, and three of the choices
bear directly on the floors. Exposures do not change between vintages. Whether a loan
is live in a vintage is drawn independently of its size and its key. And the
adversary's reading error is a declared normal error on the log scale. The register
was rebuilt once after referees found that its loan panel gave the history adversary
a signal no real register has: every loan had been forced to be live in the first
month, and a third vanished and returned in a fixed pattern. The generator now has
originations during the panel, a monthly exit hazard of 0.04 with no return, and a
stable stock of about 27{,}000 loans.
% source: docs/panel_review/10_round4_panel.md finding 1
That correction moved the results, which is a measure of how much they depend on the
generator.

The repository states what an institution would have to supply to replace the
synthetic evidence with its own. An adapter, written inside the institution and
outside the package, would emit only aggregates: de-identified key-class counts,
held-out predicted match probabilities with their outcomes for the calibration test,
external frequency benchmarks for the transport test, and aligned attack-success
indicators for every release. % source: docs/CLAIMS_EVIDENCE.md "Required real-data extension"
No direct identifier, auxiliary register or row-level confidential data would leave
the secure environment. The demographic path of the code has been run this way on two
public files (Chapter~\ref{ch:public}). The credit path has not been run on any real
register.

Three replications would tell an institution most of what the synthetic study cannot.
The first is the reading precision: the institution has to argue, from what it knows
of the sources outsiders hold, for the error with which an outsider reads exposure,
and it should report the floor across a range of precisions, as the replication
table of the long form does, and not at one value. The second is
the composition gap on its own programme: the ratio of the pooled risk to the bound the
marginals imply, which on the simulated population ranged from 1.57 to 2.05 across
register sizes from 50{,}000 to 400{,}000 units and reached 1.86 on the public Census
extract. % source: paper/generated/numbers.tex runscalesmallratio 1.57, runscalelargeratio 2.05, runrealpooledratio 1.86
The third is what a floor keeps of the credit volume. On the synthetic register the
floor of 150 keeps 20 per cent of loans and 4 per cent of exposure; on a register
with more lender concentration and regional specialisation, both of which the referees
asked for and the generator does not yet have, the share could be very different.
% source: docs/panel_review/10_round4_panel.md (central banking, open items); paper/generated/table_robustness.tex

The book's prose has to carry this limit wherever the credit floors appear, and the
chapters that report them say so. The floors of 75, 150 and none are properties of a
synthetic register under declared parameters. The decision rule that produced them is
the transferable result.

At the Bank, the chair asked Naledi in the first committee meeting of the year which
of the book's numbers she would put in front of the board. Her answer was none of the
floors and all of the method. The floors described a register the Bank did not hold.
What she could present was the record for Tomás's request, filled with the Bank's own
counts inside the secure environment, and the range of floors it gave across the
reading precisions the reviewer had accepted as plausible.

\section{Time-varying exposures and dependent loan activity}

Proposition~\ref{prop:nohistory} explains a result that surprised the referees: an
adversary who pools six vintages without knowing the target's history gains nothing
over the first vintage, and the floor stays at 1. The proposition rests on two
conditions that a real register will not meet exactly. Each loan's published key and
exposure are the same in every vintage in which it appears, and whether a loan is
live in a vintage is independent of its key, its exposure and the adversary's
information.

Both conditions fail in ordinary credit data. Exposures amortise, are drawn down and
are restructured, so a loan's published exposure differs from month to month. An
adversary who holds a reading of the target in each vintage then compares each
published value with a different reading, and the later vintages are no longer copies
of the first. Loan activity also depends on size: large loans to large firms may
persist longer, and small credit lines may open and close more often. When activity
depends on the key or the exposure, the pattern of appearances across vintages carries
information about which candidate is the target, even for an adversary who does not
know the target's own pattern. In both cases the pooled adversary without history
would gain something, and the zero gain of the synthetic register is a lower limit.

The referees listed amortising exposure among the open items, together with the
programme analysis of monthly aggregate tables and the harm that a lender-insider
causes to other lenders' relationships. % source: docs/panel_review/10_round4_panel.md central banking open items
None of these needs new theory. Each needs a generator, or better a real register,
with the feature in it, and a rerun of the joint attack. The proposition says what to
look for: the gain from pooling without history is exactly the information that
changing values and dependent activity add, and on a register without them it is zero.

The software now carries a first model of this kind, in \texttt{privacy\_risk.dynamics}:
log exposure follows an autoregressive path that continues as a latent value while a
loan is inactive, activity follows a logistic transition that can depend on the
previous state and on exposure, and the pooled attack uses absence as well as
presence. Its static limit holds exposure constant and makes activity independent of
the past, which is the assumption of Chapter~\ref{ch:credit}; it is not that chapter's
register, because the panel has no originations or exits. Chapter~\ref{ch:change}
reports what the model shows on synthetic panels. An adversary with six vintages of
history and readings identified 0.34 of targets on a frozen panel and 0.59 on one
whose exposures drift and whose activity persists in the extension's run (the
repository's rerun with seed 2026 gives 0.36 and 0.62), so the frozen register
understated this adversary's risk. The model is still code and tests, not evidence
about a real register: no floor in this book has been recomputed under it, and its
parameters have to be estimated on a governed register.

\section{Attack coverage}

Every certificate in this book covers the attacks that were declared. The admissible
release is one whose declared attacks succeed at a rate below the tolerance with stated
confidence, and the method says nothing about attacks outside the declared set. Two
documented cases show what an undeclared attack looks like. \citet{narayanan2008}
showed that subscribers in a film-rating data set released for a prediction contest
could be identified from a handful of ratings with approximate dates, using a public
rating site as the auxiliary source; the release had been reviewed without an
adversary who held such a source. The United States Census Bureau's experiments on
the tables of the 2010 census showed that published tables could be inverted to
reconstruct records \citep{garfinkel2019}, and the Bureau adopted a formal privacy
mechanism for 2020 \citep{abowd2022}. In each case the attack was not new in
principle. It was new to the declaration. Sampling gives less cover than it seems to as
well: \citet{rocher2019} showed that a model fitted to an incomplete data set can
estimate how likely a given match is to be correct.

The declared attack can also be weaker than the declared knowledge allows. The
referees found that the history adversary, who knows when the target's loan is live,
also knows when it is absent, and that the attack first used only half of that
knowledge. Using both halves raised the risk at the old floor of 40 from 0.105 to
0.179. % source: docs/panel_review/10_round4_panel.md finding 2
The attack now requires a candidate's pattern of appearances to match the target's in
both directions, but it does so through a large penalty in the score and not through
exact exclusion, and whether the maximum-posterior rule is the best attack for each
declared information set has not been proved. % source: docs/panel_review/10_round4_panel.md (differential privacy, open)
Simulation can show what a declared rule achieves. It cannot deliver the supremum over
rules that a formal guarantee delivers by construction. Audits of differentially
private training meet the same limit: an attack run in an audit gives a lower bound
on what an adversary can do, never an upper one \citep{jagielski2020,steinke2023}.

The insider is the clearest case of what a floor cannot do, and
Proposition~\ref{c18:prop-insider} states it without simulation.

\begin{proposition}[An exact insider is not governed by a floor]\label{c18:prop-insider}
Suppose a release publishes, for each record it keeps, a continuous attribute whose
values are distinct across records with probability one, and that the protection
mechanism removes records, by any rule including a minimum cell count, without
altering the published values of the records it keeps. An adversary who knows the
target's published value exactly identifies the target with probability one whenever
the target is released. Its membership-conditional risk is therefore one at every
setting of the mechanism, and its unconditional risk equals the probability that the
target is released.
\end{proposition}

\begin{example}[Pieter and an outsider on one key class]\label{c18:ex-pieter}
The numbers are illustrative. A key class at the floor keeps five loans whose published
exposures are $1.20$, $1.35$, $2.10$, $2.40$ and $3.05$ million, and the target is the loan of $2.10$.
An outsider reads the target's exposure with a log error of standard deviation $0.6$, and
suppose the reading comes out half a standard deviation high, at $2.84$. The outsider's
posterior over the five loans is $0.098$, $0.127$, $0.241$, $0.263$ and $0.271$, so it names the loan of
$3.05$ and is wrong. Pieter reads the same exposure with a log error of $0.01$. Half a standard
deviation high gives $2.11$, and his posterior puts $1.000$ on the loan of $2.10$. A higher floor can
only drop this key class from the release. If the class stays, the value $2.10$ stays with it,
and Pieter finds it. (The posteriors are the normal-error posterior on the log scale with a
flat prior over the five loans, recomputed here.)
\end{example}

The proposition explains the insider row of the decision record. An insider who reads
exposure with a log error of 0.01 is close to exact, and on the synthetic register its
success falls only from 0.772 with no suppression to 0.416 at a floor of 300.
% source: paper/bis_ifc_full.tex tab:record (insider row)
What lowers it is the removal of records, not their protection. Against such an
adversary the levers are the access route, the precision of what is published, or a
formal mechanism. Rounding exposure to one significant figure halves the insider's
success on the synthetic register. % source: docs/panel_review/10_round4_panel.md ("It halves the insider's success")

The public-file experiments add a measured view of coverage. On thirty state samples
of the Census extract, a search guided by a neural network confirmed 31.57 of 32
high-risk settings on average, and a simple domain heuristic confirmed all 32. A
quasi-random Sobol search confirmed only 22.03, but it covered 13.83 of 24 scenario
cells, against 10.07 for the neural search and 11.97 for the heuristic.
% source: paper/generated/table_public_search.tex
A search that targets what it already believes is risky finds risky settings and
covers less of the space. The cells are fixed bins of declared configuration choices,
not new attack families, so even the widest search here explores a space the analyst
wrote down. No method in the repository proposes an attack outside the declared space
and then measures it.

When Grace asked what the Bank's records could not show, Naledi's list began with
attacks. The record showed what each declared adversary achieved and how confident
the Bank could be about it. It could not show that the declared adversaries were the
ones that would come, and she wrote that sentence into the record's limits section.

\section{Eliciting the harm weight}

Chapter~\ref{ch:tolerance} showed that a stated weight on harm relative to lost
utility implies a tolerance, through the point of the risk--utility frontier that
maximises utility minus the weighted risk bound. The result is useful only if a
committee can state the weight, and no committee currently does. The trade has been
framed before: \citet{abowd2019} treat the balance of privacy loss against accuracy as
a social choice with a rate of exchange between the two, and \citet{hotz2022} argue
that statistical agencies should weigh the usefulness of their data beside privacy
when they choose a protection. The repository
reports the implied tolerance at six weights on the simulated population: 0.464 at a
weight of 0.25, 0.345 at 0.5, 0.145 at 1 and at 2, and 0.081 at 5 and at 10.
% source: paper/generated/table_tolerance.tex; paper/generated/numbers.tex runtolerancekappa*

The figures contain a practical point. Because the frontier is a finite set of
settings, the implied tolerance is a step function of the weight: every weight between
1 and 2 implies the same setting, and so does every weight between 5 and 10. A
committee does not need to state its weight to two decimal places. It needs to place
the weight in an interval wide enough to contain the steps, and it can do so by
choosing between pairs of settings shown with what each keeps and what each risks.
Whether committees can make such choices consistently, and whether their past release
decisions identify a weight, are empirical questions. The long form proposes estimating
the weight from revealed institutional behaviour, since agencies have made these
decisions for decades and their past choices constrain it under the framework. That
study has not been done.

\begin{example}[Placing the chair's weight]\label{c18:ex-weight}
The numbers are illustrative. A frontier has four settings, with upper risk bounds of
$0.40$, $0.25$, $0.12$ and $0.05$ and utilities of $1.00$, $0.90$, $0.78$ and $0.60$. The setting chosen is
the one that maximises utility minus the weight times the bound. The second setting beats
the first when the weight exceeds $0.10/0.15=0.667$, the third beats the second above
$0.12/0.13=0.923$, and the fourth beats the third above $0.18/0.07=2.571$. Every weight between
$0.923$ and $2.571$ therefore implies the tolerance of $0.12$. The chair need not say whether her
weight is 1 or 2. She has to answer two questions: would she give up $0.12$ of utility to
lower the bound from $0.25$ to $0.12$, and would she give up a further $0.18$ to lower it to
$0.05$? A yes and a no place her weight in the interval, and the tolerance follows.
\end{example}

The harm weight also hides a scope limit. The framework represents harm by a scalar
within a declared harm class. Harms that fall on groups and not on individuals, and
harms that arise from correct inference about a population and not from the
identification of a person, are outside it \citep{skinner2012,dwork2014}.

The chair had never been asked for a weight on harm, and when Naledi asked she gave a
tolerance instead. Naledi wrote both into the record for Tom\'as's request: the
tolerance the chair signed, and a note that the weight it implies had not been
elicited. Grace later put that note first on the Bank's list of open items.

\section{The model-proposed attack line}

Chapter~\ref{ch:proposals} reported a pipeline in which a model proposes scenarios and
bets, a frozen plan records them, and an independent audit decides. Its results are
mostly negative, and the remaining questions are best stated with those results
beside them. On the synthetic benchmark, trained to convergence over ten seeds, the
logistic and deep proposal models reach held-out losses that cannot be told apart.
% source: README.md (SENTRY-R section)
On the public Census file, one- and three-layer networks improved on plain logistic
regression in three overlapping state splits, and a logistic model with explicit
pairwise products and squares did better than both.
% source: docs/SIX_PERSPECTIVE_REVIEW.md findings
The evidence supports a role for nonlinearity, not for depth.

The language-model arm is where the gap between what is built and what is measured is
widest. A hosted pilot was prepared, with prompts saved and hashed before collection,
and its first request returned an error with no proposal; collection stopped without
retrying, and no hosted model response exists.
% source: docs/LLM_PILOT_STATUS_2026_10_23.md
A second run used a small open-weight model with four billion parameters on the
author's own machine. It produced a valid list of proposals for 20 of 30 states within
three attempts, and 50 of its 70 attempts were rejected, most for repeated or
off-catalogue proposals. On those 20 states its proposals covered 6.05 scenario cells
on average, against 13.65 for the Sobol search, and confirmed 27.60 of 32 settings. In
19 of the 20 states every proposal used the finest geography, in the catalogue's own
order. % source: paper/generated/numbers_llm_local.json
The model copied the catalogue. Its run settles that this small model, under this
prompt, adds no coverage over quasi-random or neural search. It does not settle what a
larger model, a prompt that names the coverage goal, or a worked example would do, and
any of those is a new experiment to be declared before it is run.

There is therefore no genuine comparison, at a fixed evaluation budget, of a
capable language model against the strongest baseline that uses no language model, which is the
primary contrast of the confirmatory protocol in the repository.
% source: EXPERIMENT_PROTOCOL.md
The infrastructure for it exists: a replay contract that binds each response to the
hash of the exact prompt, a schema that refuses counts and code, and a ledger that
records the order of proposal, review, measurement, training, freeze and audit. Until
the comparison is run, the claim the line can support is architectural. A model can
propose scenarios into an audit without being able to change its error statement, and
a poor model costs attacks, not validity.

\section{Claims the method cannot make}

Some limits are not gaps that more work will close. They follow from what the method
is, and a release record should state them.

The method cannot establish legal compliance. Data protection law asks whether
identification is reasonably likely by the means an adversary could use, and gives no
number \citep{gdpr2016}. The statistical confidentiality duty of the European System of
Central Banks likewise names no probability \citep{ec2533}. A tolerance is the
institution's written reading of that phrase for a declared adversary and harm. The
method makes the reading explicit, owned and dated, and leaves to lawyers the question
whether it is a reasonable reading.

The method cannot protect against adversaries who were not declared. Its certificate
covers registered attacks. An adversary with better auxiliary information, a new
linkage method or a different notion of success lies outside the claim, and
Proposition~\ref{c18:prop-insider} shows that one such adversary, an insider who knows a
published value exactly, defeats every floor. The falsification tests and the expiry
date narrow this gap; they do not close it. The method gives no differential privacy
guarantee either, and where a guarantee over all attacks is required, a formal
mechanism with an accountant is the tool (Chapter~\ref{ch:dp}).

Naledi's limits section for Tom\'as's record grew with each of these paragraphs. It
names the declared adversaries and says the certificate covers them alone. It names
Pieter and says that no floor governs him. And it leaves to the Bank's lawyers the
question whether a tolerance of $0.12$ is a reasonable reading of the law.

The method cannot enforce its own controls. The implementation is a library and not a
security boundary, the registry service that a deployment needs is specified and not
built beyond a single-host ledger (Chapter~\ref{ch:planner}), and the hashes a run writes detect change but do not prove who made it.
% source: docs/DEPLOYMENT_AND_PENTEST_2026_10_23.md; paper/bis_ifc_full.tex sec:security
Separation of duties, external time stamps and audit inside the secure environment are
institutional controls, of the kind that guidance on de-identifying government data sets places under governance \citep{garfinkel2023}, and the method depends on them.

The method also cannot certify what it has not shown in the open. The papers describe
the pipeline as open, and a reader's ability to rerun a number is part of the
evidence for it. At the time of writing the repository is private, and either it is
made public or that description changes before the papers or this book circulate.
% source: book/REVIEW.md point 5

\section{Open statistical questions}

Four extensions look tractable, and the long form names them. The selected point of
the risk--utility frontier is chosen by maximising an estimate, so it is biased in the
way any such choice is, and the frontier needs a proper inferential treatment. The
risk budget, which allocates a programme tolerance across releases, controls only the
union event; a joint accountant that admits dependence between releases would need a
composition theorem for scenario-conditional risk under a stated class of dependence.
The implied tolerance could be estimated from revealed behaviour, as above. And the
calibration test of Chapter~\ref{ch:falsification} could be developed into a formal
specification test with a distribution under the null, which would replace the
analyst's band with a level. The sequential audit has open items of its own: a
truncated Beta mixture in place of the capped bet, exact group-sequential boundaries,
and a bound on the expected length of the rule actually used.

Tom\'as read this list when the committee sent him its decision. Two of the items, the
inference for the frontier and the joint accountant, lay close to his own field, and he
asked whether the Bank would share the synthetic register so that he could work on them.
Naledi added the request to the inventory as a fourth product, with a row of its own.
% source: docs/panel_review/10_round4_panel.md (sequential statistics, open)

\section{How an institution should turn a measured risk into a release decision}

Chapter~\ref{ch:decision} asked one question: what must an institution add to a measured
risk before the result can be called a decision? Its brief answer named three additions: the
attack the figure describes, an upper bound in place of the estimate, and a tolerance with an
owner. The chapters since have added three more: the programme as the object assessed, a test
of the model before its figure is used, and a fixed order of work kept in a record. Part VII
added a condition on learning: what the institution learns about attacks may plan its audits,
but it may not stand in for them. The answer, in the order an institution should act on it, is
as follows.

An institution should first say what it is measuring. Disclosure risk is the
probability that a declared adversary, holding declared information, names the right
record for a declared set of targets. A figure without that declaration has no
referent, and the familiar measures are answers to different declarations
(Chapters~\ref{ch:scenario} and~\ref{ch:measures}).

It should then treat the figure as an estimate. The decision compares an exact upper
confidence bound with the tolerance, not a point estimate, because a point estimate at
the tolerance passes about half the time. The number of simulated attacks is fixed in
advance from the tolerance and the error level, a protection setting is chosen by
testing a fixed sequence from the most protective down and confirmed on fresh
targets, and an audit that is watched as it runs uses a sequence that stays valid
under watching (Chapters~\ref{ch:estimate} to~\ref{ch:monitoring}).

The tolerance has to belong to someone. It is written before the risk is measured,
with a named owner, an expiry date, a harm class, an error level, the legal basis and
the parties it protects, the access route and, for a programme, a horizon. It can be
stated directly or read off the risk--utility frontier from a stated weight on harm
(Chapters~\ref{ch:thresholds} and~\ref{ch:tolerance}).

The object assessed is the programme, not the release. When records keep a pseudonym
across releases, no function of the per-release risks bounds an adversary who pools
them, so the joint attack is simulated directly and compared with a programme
tolerance over the declared horizon (Chapter~\ref{ch:programmes}).

The model behind the figure is tested before the figure is used, against bands fixed
in advance, and a failed test makes the assessment invalid. The figure is not adjusted
(Chapter~\ref{ch:falsification}). The whole sequence is carried out in a fixed order
by separate people, recorded in a registry that holds every audit, and filed as a
decision record with a row for every harm, including those for which no tolerance has
been set (Chapter~\ref{ch:protocol}).

What the institution learns over time changes how it plans, not what it certifies. A
system that keeps searching for weaknesses may propose attacks, size audits and set their
bets, but each attempt is charged to an error budget before its first outcome, and the
release is certified only by fresh attacks through the same gate (Chapters~\ref{ch:learning}
and~\ref{ch:planner}). A register that changes over time gives a watching adversary more to
work with, so a frozen register cannot be assumed to be the safe case
(Chapter~\ref{ch:change}). A model of how the population moves may enter a decision only
after it forecasts held-out periods better than a simpler baseline, which no model in this
book has yet done (Chapter~\ref{ch:dynamics}). And when the adversary changes, the
institution measures how sensitive its decision is and reassesses; a new kind of computer
changes the encryption that protects the register, and leaves unchanged what a release
shows (Chapter~\ref{ch:chaos}).

Tom\'as's request went through each of these steps. Its record names the adversaries, sets
upper bounds against a tolerance the chair signed, assesses his three products as one
programme, files the model's tests against bands fixed in advance and names the remedy the
committee chose, the secure room. Grace could check every entry against its date. None of
this made the floors less uncertain. It made the decision one that someone else could
follow and dispute.

The release goes out when the upper bound for every declared scenario, pooled ones
included, lies at or below its tolerance, the model has passed its tests, and the
owner has signed. Otherwise the record names the remedy, and the changed product is
assessed again.

That rule is reproducible: two analysts with the same declarations and the same
simulation output reach the same decision, and an auditor can check afterwards that
the decision followed from the evidence and that the evidence was gathered in the
stated order. The rule does not remove judgement. It places the judgement in the
declarations and the tolerance, where it can be seen, disagreed with and revised, and
it leaves the probability calculation open to being proved wrong. What remains open,
as this chapter has set out, is the evidence from real registers that would show how
large the numbers are, and the declarations that no method can make on an
institution's behalf.

\begin{inpractice}[In practice: stating the limits in the record]
\setlength{\emergencystretch}{3em}%
Add a limits section to every decision record. State that the certificate covers the
declared attacks only, name any adversary that was considered and not declared, and
say why. Report floors across a range of the parameters that decide them, the
adversary's reading precision above all, and not at one value. State whether
exposures change between vintages and whether loan activity depends on size, since
the no-history result needs both to fail. Name any insider separately and govern it
through access. Record that the tolerance is the institution's reading of its legal
duty, not a finding of compliance.

\smallskip\noindent\textit{In the software.} The floors across reading precisions and
register draws come from \texttt{scripts/\allowbreak{}programme\_\allowbreak{}robustness.\allowbreak{}py} and
\texttt{scripts/\allowbreak{}robustness\_\allowbreak{}table.\allowbreak{}py}. The public-file search and the language-model
arm are \texttt{privacy\_\allowbreak{}risk.\allowbreak{}injector\_\allowbreak{}search},
\texttt{privacy\_\allowbreak{}risk.\allowbreak{}llm\_\allowbreak{}pilot\_\allowbreak{}local} and
\texttt{scripts/\allowbreak{}public\_\allowbreak{}injector\_\allowbreak{}search.\allowbreak{}py}. The adapter for a real register is not
in the package; \texttt{privacy\_\allowbreak{}risk.\allowbreak{}realdata} is the pattern for writing one.
\end{inpractice}

\section{Proofs}

\begin{proof}[Proof of Proposition~\ref{c18:prop-insider}]
Let $v_j$ be the published value of the continuous attribute for each released record
$j$, and let $t$ be a released target with value $v_t$. The mechanism keeps $v_t$
unaltered, so the adversary, who knows $v_t$ exactly, can form the set
$\{j : v_j = v_t\}$ of released records sharing it. With probability one the values
are distinct across records, so the set is $\{t\}$, and naming its only member
identifies the target. The membership-conditional risk is therefore one for every rule
that removes records, at every setting. The unconditional risk is the product of the
probability that the target is released and the membership-conditional risk, by the
factorisation of Chapter~\ref{ch:scenario}, and so equals the probability of release.
\end{proof}
